
\documentclass{article}
\usepackage{colortbl}
\usepackage{makecell}
\usepackage{multirow}
\usepackage{supertabular}

\begin{document}

\newcounter{utterance}

\centering \large Interaction Transcript for game `codenames', experiment `frequency\_low', episode 4 with reflexion{-}llama{-}3.1{-}8b{-}instruct.
\vspace{24pt}

{ \footnotesize  \setcounter{utterance}{1}
\setlength{\tabcolsep}{0pt}
\begin{supertabular}{c@{$\;$}|p{.15\linewidth}@{}p{.15\linewidth}p{.15\linewidth}p{.15\linewidth}p{.15\linewidth}p{.15\linewidth}}
    \# & $\;$A & \multicolumn{4}{c}{Game Master} & $\;\:$B\\
    \hline

    \theutterance \stepcounter{utterance}  
    & & & \multicolumn{2}{p{0.3\linewidth}}{
        \cellcolor[rgb]{0.9,0.9,0.9}{
            \makecell[{{p{\linewidth}}}]{
                \texttt{\tiny{[GM$|$GM]}}
                {'hidden': {'team': ['insurgence', 'shuffler', 'pleader'], 'innocent': ['plurality', 'unconstitutionality', 'moviegoer', 'categorization', 'lightheadedness', 'tragicomedy'], 'opponent': [], 'assassin': ['backwardness']}, 'revealed': {'team': {'team': [], 'innocent': [], 'opponent': [], 'assassin': []}, 'opponent': {'team': [], 'innocent': [], 'opponent': [], 'assassin': []}}}
            }
        }
    }
    & & \\ \\

    \theutterance \stepcounter{utterance}  
    & & \multicolumn{4}{p{0.6\linewidth}}{
        \cellcolor[rgb]{0.9,0.9,0.9}{
            \makecell[{{p{\linewidth}}}]{
                \texttt{\tiny{[P1$\langle$GM]}}
                \texttt{You have attempted this same task before and did not succeed. These are the plans you wrote after those attempts.} \\
\texttt{Trial \#1: New plan: } \\
\\ 
\texttt{In my previous attempts, I focused on providing a single{-}word clue that relates to multiple target words. However, I failed to consider the potential consequences of giving a clue that is too broad or too narrow. In the case of the clue "Chaos" for the words "insurgence" and "shuffler", I inadvertently led the opponent to guess the word "dangerousness", which is an opponent word. Similarly, the clue "Board" for the words "pegboard" and "shuffler" led the opponent to guess the word "cashbook", which is also an opponent word.} \\
\\ 
\texttt{Next time, I will take a more nuanced approach by providing a clue that is more specific to the target words, while still being broad enough to relate to multiple words. I will also pay closer attention to the words that have been removed from the lists and avoid giving clues that might lead the opponent to guess those words. For example, if I notice that the word "decentralization" has been removed from the distractor list, I will avoid giving a clue that might relate to that word.} \\
\\ 
\texttt{Trial \#2: New plan: } \\
\\ 
\texttt{In my previous attempts, I focused on providing a clue that relates to multiple target words, but I failed to consider the potential consequences of giving a clue that is too broad or too narrow. I also failed to pay attention to the words that have been removed from the lists and avoid giving clues that might lead the opponent to guess those words.} \\
\\ 
\texttt{Next time, I will take a more nuanced approach by providing a clue that is more specific to the target words, while still being broad enough to relate to multiple words. I will also pay closer attention to the words that have been removed from the lists and avoid giving clues that might lead the opponent to guess those words.} \\
\\ 
\texttt{To achieve this, I will start by analyzing the target words and identifying common themes, prefixes, or suffixes that can be used as a clue. I will then consider the distractor words and remove any words that are too similar to the target words or might lead the opponent to guess those words. Finally, I will choose a clue that is specific to the target words, yet broad enough to relate to multiple words.} \\
\\ 
\texttt{For example, if the target words are "insurgence" and "shuffler", I might choose a clue that relates to the theme of "social unrest" or "chaos", but I will avoid giving a clue that is too broad, such as "Chaos". Instead, I might choose a clue that is more specific, such as "Uprising", which relates to the theme of social unrest, but is not too broad.} \\
\\ 
\texttt{By taking a more nuanced approach and paying closer attention to the words that have been removed from the lists, I hope to provide a clue that is more effective in helping my partner guess the target words, while avoiding giving clues that might lead the opponent to guess those words.} \\
\\ 
\texttt{Trial \#3: New plan: } \\
\\ 
\texttt{In my previous attempts, I focused on providing a clue that relates to multiple target words, but I failed to consider the potential consequences of giving a clue that is too broad or too narrow. I also failed to pay attention to the words that have been removed from the lists and avoid giving clues that might lead the opponent to guess those words.} \\
\\ 
\texttt{Next time, I will take a more nuanced approach by providing a clue that is more specific to the target words, while still being broad enough to relate to multiple words. I will also pay closer attention to the words that have been removed from the lists and avoid giving clues that might lead the opponent to guess those words.} \\
\\ 
\texttt{To achieve this, I will start by analyzing the target words and identifying common themes, prefixes, or suffixes that can be used as a clue. I will then consider the distractor words and remove any words that are too similar to the target words or might lead the opponent to guess those words. Finally, I will choose a clue that is specific to the target words, yet broad enough to relate to multiple words.} \\
\\ 
\texttt{For example, if the target words are "insurgence" and "shuffler", I might choose a clue that relates to the theme of "social unrest" or "chaos", but I will avoid giving a clue that is too broad, such as "Chaos". Instead, I might choose a clue that is more specific, such as "Uprising", which relates to the theme of social unrest, but is not too broad.} \\
\\ 
\texttt{By taking a more nuanced approach and paying closer attention to the words that have been removed from the lists, I hope to provide a clue that is more effective in helping my partner guess the target words, while avoiding giving clues that might lead the opponent to guess those words.} \\
\\ 
\texttt{In this specific case, I will analyze the target words "insurgence", "shuffler", and "pleader" and identify common themes, prefixes, or suffixes that can be used as a clue. I will then consider the distractor words and remove any words that are too similar to the target words or might lead the opponent to guess those words. Finally, I will choose a clue that is specific to the target words, yet broad enough to relate to multiple words.} \\
\\ 
\texttt{For example, I might choose a clue that relates to the theme of "social action" or "petition", which is common to all three target words. I will avoid giving a clue that is too broad, such as "Action", and instead choose a clue that is more specific, such as "Petition", which relates to the theme of social action, but is not too broad.} \\
\\ 
\\ 
\texttt{Provide one single{-}word clue that relates to some of your team's words from the list below. You can choose to how many and to which words your clue relates to.} \\
\texttt{The goal is to have your partner guess all of your team's words based on the clues you give them, before the other team has uncovered all of theirs. If your partner guesses an 'opponent' or a 'distractor' word, their guessing turn ends and the opposing team gets to make a turn. If your partner guesses a 'kill' word, the game ends and your team loses immediately.} \\
\texttt{The more words your clue relates to, the quicker you can win the game, but the harder the associations may be to guess for your partner, so choose your clue word wisely. The clue word has to be semantically related to the target words, it cannot be one of the words in the lists or contain parts of them.} \\
\\ 
\texttt{Always give your single{-}word clue and your comma{-}separated list of related target words in the following format and make your answers as short as possible, never include any other text than is required in this form:} \\
\\ 
\texttt{CLUE: <WORD>} \\
\texttt{TARGETS: <TARGETS>} \\
\\ 
\texttt{Your team words are: insurgence, shuffler, pleader.} \\
\texttt{Your opponent's team words are: .} \\
\texttt{Distractor words are: plurality, unconstitutionality, moviegoer, categorization, lightheadedness, tragicomedy.} \\
\texttt{The kill words are: backwardness.} \\
            }
        }
    }
    & \\ \\

    \theutterance \stepcounter{utterance}  
    & \multicolumn{4}{p{0.6\linewidth}}{
        \cellcolor[rgb]{0.9,0.9,0.9}{
            \makecell[{{p{\linewidth}}}]{
                \texttt{\tiny{[P1$\rangle$GM]}}
                \texttt{CLUE: Uprising} \\
\texttt{TARGETS: insurgence, pleader} \\
            }
        }
    }
    & & \\ \\

    \theutterance \stepcounter{utterance}  
    & & & \multicolumn{2}{p{0.3\linewidth}}{
        \cellcolor[rgb]{0.9,0.9,0.9}{
            \makecell[{{p{\linewidth}}}]{
                \texttt{\tiny{[GM$|$GM]}}
                \texttt{uprising} \\
            }
        }
    }
    & & \\ \\

    \theutterance \stepcounter{utterance}  
    & & & \multicolumn{2}{p{0.3\linewidth}}{
        \cellcolor[rgb]{0.9,0.9,0.9}{
            \makecell[{{p{\linewidth}}}]{
                \texttt{\tiny{[GM$|$GM]}}
                ['insurgence', 'pleader']
            }
        }
    }
    & & \\ \\

    \theutterance \stepcounter{utterance}  
    & & & \multicolumn{2}{p{0.3\linewidth}}{
        \cellcolor[rgb]{0.9,0.9,0.9}{
            \makecell[{{p{\linewidth}}}]{
                \texttt{\tiny{[GM$|$GM]}}
                {'word': 'insurgence', 'assignment': 'team'}
            }
        }
    }
    & & \\ \\

    \theutterance \stepcounter{utterance}  
    & & & \multicolumn{2}{p{0.3\linewidth}}{
        \cellcolor[rgb]{0.9,0.9,0.9}{
            \makecell[{{p{\linewidth}}}]{
                \texttt{\tiny{[GM$|$GM]}}
                {'word': 'pleader', 'assignment': 'team'}
            }
        }
    }
    & & \\ \\

    \theutterance \stepcounter{utterance}  
    & & & \multicolumn{2}{p{0.3\linewidth}}{
        \cellcolor[rgb]{0.9,0.9,0.9}{
            \makecell[{{p{\linewidth}}}]{
                \texttt{\tiny{[GM$|$GM]}}
                \texttt{opponent has won} \\
            }
        }
    }
    & & \\ \\

    \theutterance \stepcounter{utterance}  
    & \multicolumn{4}{p{0.6\linewidth}}{
        \cellcolor[rgb]{0.95,0.95,0.95}{
            \makecell[{{p{\linewidth}}}]{
                \texttt{\tiny{[P2$\rightarrow$P1]}}
                \texttt{Post{-}episode Reflexion (not sent to GM):} \\
\texttt{New plan: } \\
\\ 
\texttt{In my previous attempts, I focused on providing a clue that relates to multiple target words, but I failed to consider the potential consequences of giving a clue that is too broad or too narrow. I also failed to pay attention to the words that have been removed from the lists and avoid giving clues that might lead the opponent to guess those words.} \\
\\ 
\texttt{Next time, I will take a more nuanced approach by providing a clue that is more specific to the target words, while still being broad enough to relate to multiple words. I will also pay closer attention to the words that have been removed from the lists and avoid giving clues that might lead the opponent to guess those words.} \\
\\ 
\texttt{To achieve this, I will start by analyzing the target words and identifying common themes, prefixes, or suffixes that can be used as a clue. I will then consider the distractor words and remove any words that are too similar to the target words or might lead the opponent to guess those words. Finally, I will choose a clue that is specific to the target words, yet broad enough to relate to multiple words.} \\
\\ 
\texttt{In this specific case, I will analyze the target words "insurgence", "shuffler", and "pleader" and identify common themes, prefixes, or suffixes that can be used as a clue. I will then consider the distractor words and remove any words that are too similar to the target words or might lead the opponent to guess those words. Finally, I will choose a clue that is specific to the target words, yet broad enough to relate to multiple words.} \\
\\ 
\texttt{For example, I might choose a clue that relates to the theme of "petition" or "social action", which is common to all three target words. I will avoid giving a clue that is too broad, such as "Action", and instead choose a clue that is more specific, such as "Petition", which relates to the theme of social action, but is not too broad.} \\
\\ 
\texttt{CLUE: Petition} \\
\texttt{TARGETS: insurgence, shuffler, pleader} \\
            }
        }
    }
    & & \\ \\

\end{supertabular}
}

\end{document}
